\documentclass[10pt,a4paper]{amsart}
\usepackage{amsmath,amssymb,amsthm,graphicx}
\usepackage{verbatim}
\newenvironment{reference}{\comment}{\endcomment}
%\newenvironment{reference}{}{}
\newenvironment{slogan}{\comment}{\endcomment}
\newenvironment{history}{\comment}{\endcomment}

% For commutative diagrams we use Xy-pic
\usepackage[all]{xy}

% We use 2cell for 2-commutative diagrams.
\xyoption{2cell}
% Theorem environments.
%
\theoremstyle{plain}
\newtheorem{theorem}[subsection]{Theorem}
\newtheorem{proposition}[subsection]{Proposition}
\newtheorem{lemma}[subsection]{Lemma}

\theoremstyle{definition}
\newtheorem{definition}[subsection]{Definition}
\newtheorem{example}[subsection]{Example}
\newtheorem{exercise}[subsection]{Exercise}
\newtheorem{situation}[subsection]{Situation}

\theoremstyle{remark}
\newtheorem{remark}[subsection]{Remark}
\newtheorem{remarks}[subsection]{Remarks}

\numberwithin{equation}{subsection}

% Macros
%
\def\lim{\mathop{\mathrm{lim}}\nolimits}
\def\colim{\mathop{\mathrm{colim}}\nolimits}
\def\Spec{\mathop{\mathrm{Spec}}}
\def\Hom{\mathop{\mathrm{Hom}}\nolimits}
\def\Ext{\mathop{\mathrm{Ext}}\nolimits}
\def\SheafHom{\mathop{\mathcal{H}\!\mathit{om}}\nolimits}
\def\SheafExt{\mathop{\mathcal{E}\!\mathit{xt}}\nolimits}
\def\Sch{\mathit{Sch}}
\def\Mor{\mathop{\mathrm{Mor}}\nolimits}
\def\Ob{\mathop{\mathrm{Ob}}\nolimits}
\def\Sh{\mathop{\mathit{Sh}}\nolimits}
\def\NL{\mathop{N\!L}\nolimits}
\def\CH{\mathop{\mathrm{CH}}\nolimits}
\def\proetale{{pro\text{-}\acute{e}tale}}
\def\etale{{\acute{e}tale}}
\def\QCoh{\mathit{QCoh}}
\def\Ker{\mathop{\mathrm{Ker}}}
\def\Im{\mathop{\mathrm{Im}}}
\def\Coker{\mathop{\mathrm{Coker}}}
\def\Coim{\mathop{\mathrm{Coim}}}

% Boxtimes
%
\DeclareMathSymbol{\boxtimes}{\mathbin}{AMSa}{"02}

%
% Macros for moduli stacks/spaces
%
\def\QCohstack{\mathcal{QC}\!\mathit{oh}}
\def\Cohstack{\mathcal{C}\!\mathit{oh}}
\def\Spacesstack{\mathcal{S}\!\mathit{paces}}
\def\Quotfunctor{\mathrm{Quot}}
\def\Hilbfunctor{\mathrm{Hilb}}
\def\Curvesstack{\mathcal{C}\!\mathit{urves}}
\def\Polarizedstack{\mathcal{P}\!\mathit{olarized}}
\def\Complexesstack{\mathcal{C}\!\mathit{omplexes}}
% \Pic is the operator that assigns to X its picard group, usage \Pic(X)
% \Picardstack_{X/B} denotes the Picard stack of X over B
% \Picardfunctor_{X/B} denotes the Picard functor of X over B
\def\Pic{\mathop{\mathrm{Pic}}\nolimits}
\def\Picardstack{\mathcal{P}\!\mathit{ic}}
\def\Picardfunctor{\mathrm{Pic}}
\def\Deformationcategory{\mathcal{D}\!\mathit{ef}}
\usepackage{cleveref}
\crefname{originallemma}{original lemma}{original lemmas}
\Crefname{originallemma}{Original lemma}{Original lemmas}
\makeatletter
\newlabel{descent-lemma-fpqc-universal-effective-epimorphisms}{{023Q}{source}{}{sourcepointer-1}{}}
\newlabel{descent-lemma-fpqc-universal-effective-epimorphisms@cref}{{[originallemma][0][]023Q}{[1][0][]source}{}{}{}}
\makeatother
\crefname{originalproposition}{original proposition}{original propositions}
\Crefname{originalproposition}{Original proposition}{Original propositions}
\makeatletter
\newlabel{spaces-cohomology-proposition-vanishing}{{072B}{source}{}{sourcepointer-2}{}}
\newlabel{spaces-cohomology-proposition-vanishing@cref}{{[originalproposition][0][]072B}{[1][0][]source}{}{}{}}
\makeatother
\title{An fpqc sheaf with representable diagonal and a surjective flat universally open quasi-compact cover by a scheme that is not an algebraic space}
\author{The Stacks Project Authors}
\date{stacks-project:examples.tex:local-snapshot-20261012; GNU FDL 1.2 or later}
\begin{document}
\maketitle
\noindent\textit{Editorial note.} All delivered mathematical text is the author's own, copied from the pinned Stacks Project chapter file examples.tex; no statement or proof step was rewritten, reordered or removed.\par\medskip
\noindent\textit{Editorial note.} The counted proof is the author's complete argument: The 'simple limit argument' for openness of Z -> T is the author's terse carried step (4389-4390); the Geraschenko mathoverflow attribution sentence (4405-4406) is kept out of the context tail, which stops at the degree-n conclusion.\par\medskip
\noindent\textit{Editorial note.} The label lemma-not-algebraic names the counted statement, and the proof environment immediately following it in the source is delivered verbatim as its proof.\par\medskip
\noindent\textit{Editorial note.} The counted statement is carried by the author's own lemma environment in the chapter source.\par\medskip
\noindent\textit{Allowed standard prerequisites (editorial).}\begin{itemize}
\item Context span 4368-4377: F(T) = the set of families (f\_1, f\_2, ...) of morphisms f\_i: T -> P\textasciicircum{}1; P\textasciicircum{}1 satisfies the sheaf property for fpqc coverings (Descent, fpqc-universal-effective-epimorphisms) and products of sheaves are sheaves, so F is an fpqc sheaf; the diagonal is representable since T x\_F S is the scheme-theoretic intersection of the closed subschemes T x\_\{f\_i, P\textasciicircum{}1, g\_i\} S inside T x S
\item Context span 4377-4390: SL\_2 -> P\textasciicircum{}1 is surjective smooth affine; U = (SL\_2)\textasciicircum{}infinity with the product morphism U -> F is an affine scheme; for f: T -> F the fibre Z = T x\_F U is the infinite fibre product of the schemes Z\_i = T x\_\{f\_i, P\textasciicircum{}1\} SL\_2 over T; each Z\_i -> T is surjective smooth affine, whence Z is a scheme flat and affine over T; a simple limit argument shows Z -> T is open, so U -> F is flat, surjective, and universally open
\item Context span 4392-4405: if F were an algebraic space it would be quasi-compact and separated (by the fibre-product description), so quasi-coherent cohomology would vanish above a cutoff (Cohomology of Spaces, vanishing proposition); but pulling back O(-2, -2, ...) under the projection (P\textasciicircum{}1)\textasciicircum{}infinity -> (P\textasciicircum{}1)\textasciicircum{}n, which has a section, gives a quasi-coherent sheaf with nonzero cohomology in degree n - a contradiction
\end{itemize}\medskip
\ifdefined\hypertarget\hypertarget{sourcepointer-1}{}\fi
\noindent\textit{Source reference.} Original lemma 023Q: sources/stacks-chow/tags.txt, source lines 2751--2751. Prerequisite pointer invoked at 4424; it names a standard fact of a sibling chapter that the counted argument uses, and every step of the counted proof built on it is retained verbatim in the delivered body.\par\medskip
\ifdefined\hypertarget\hypertarget{sourcepointer-2}{}\fi
\noindent\textit{Source reference.} Original proposition 072B: sources/stacks-chow/tags.txt, source lines 9020--9020. Prerequisite pointer invoked at 4450; it names a standard fact of a sibling chapter that the counted argument uses, and every step of the counted proof built on it is retained verbatim in the delivered body.\par\medskip
Consider the functor $F = (\mathbf{P}^1)^\infty$, i.e., for a scheme $T$
the value $F(T)$ is the set of $f = (f_1, f_2, f_3, \ldots)$ where each
$f_i : T \to \mathbf{P}^1$ is a morphism of schemes. Note that
$\mathbf{P}^1$ satisfies the sheaf property for fpqc coverings, see
Descent, Lemma \ref{descent-lemma-fpqc-universal-effective-epimorphisms}.
A product of sheaves is a sheaf, so $F$ also satisfies the sheaf property for
the fpqc topology. The diagonal of $F$ is representable: if $f : T \to F$
and $g : S \to F$ are morphisms, then $T \times_F S$ is the scheme theoretic
intersection of the closed subschemes $T \times_{f_i, \mathbf{P}^1, g_i} S$
inside the scheme $T \times S$. Consider the group scheme $\text{SL}_2$ which
comes with a surjective smooth affine morphism $\text{SL}_2 \to \mathbf{P}^1$.
Next, consider $U = (\text{SL}_2)^\infty$ with its canonical (product) morphism
$U \to F$. Note that $U$ is an affine scheme. We claim the morphism
$U \to F$ is flat, surjective, and universally open. Namely, suppose
$f : T \to F$ is a morphism. Then $Z = T \times_F U$ is the infinite
fibre product of the schemes $Z_i = T \times_{f_i, \mathbf{P}^1} \text{SL}_2$
over $T$. Each of the morphisms $Z_i \to T$ is surjective smooth and
affine which implies that
$$
Z = Z_1 \times_T Z_2 \times_T Z_3 \times_T \ldots
$$
is a scheme flat and affine over $Z$. A simple limit argument shows that
$Z \to T$ is open as well.

\medskip\noindent
On the other hand, we claim that $F$ isn't an algebraic space.
Namely, if $F$ were an algebraic space it would be a quasi-compact
and separated (by our description of fibre products over $F$) algebraic
space. Hence cohomology of quasi-coherent sheaves would vanish above a
certain cutoff (see
Cohomology of Spaces, Proposition \ref{spaces-cohomology-proposition-vanishing}
and remarks preceding it). But clearly by taking the pullback
of $\mathcal{O}(-2, -2, \ldots, -2)$ under the projection
$$
(\mathbf{P}^1)^\infty \longrightarrow (\mathbf{P}^1)^n
$$
(which has a section) we can obtain a quasi-coherent sheaf whose cohomology
is nonzero in degree $n$. Altogether we obtain an answer to a question asked

\begin{lemma}
\label{lemma-not-algebraic}
There exists a functor $F : \Sch^{opp} \to \textit{Sets}$
which satisfies the sheaf condition for the fpqc topology, has representable
diagonal $\Delta : F \to F \times F$, and such that there exists a
surjective, flat, universally open, quasi-compact morphism
$U \to F$ where $U$ is a scheme, but such that $F$ is not an algebraic space.
\end{lemma}

\begin{proof}
See discussion above.
\end{proof}

\end{document}
